% !TITLE 卖炭翁
% !AUTHOR 白居易
% !DYNASTY 唐
% !GENRE 七言古诗
% !ORDER 46

\begin{poem}
卖炭翁\textsuperscript{1}，伐薪\textsuperscript{2}烧炭南山中。
满面尘灰烟火色，两鬓苍苍十指黑\textsuperscript{3}。
卖炭得钱何所营\textsuperscript{4}？身上衣裳口中食。
可怜身上衣正单，心忧炭贱愿天寒\textsuperscript{5}。
夜来城外一尺雪，晓驾炭车辗冰辙\textsuperscript{6}。
牛困人饥日已高，市南门外泥中歇。
翩翩两骑来是谁？黄衣使者白衫儿\textsuperscript{7}。
手把文书口称敕\textsuperscript{8}，回车叱牛牵向北\textsuperscript{9}。
一车炭，千余斤，宫使驱将惜不得\textsuperscript{10}。
半匹红纱一丈绫\textsuperscript{11}，系向牛头充炭直\textsuperscript{12}。
\end{poem}

\begin{annotations}
\notetext{1}{卖炭翁：一个卖炭的老人。}
\notetext{2}{伐薪：砍柴。薪（xīn），柴火。在终南山里砍柴烧炭。}
\notetext{3}{苍苍十指黑：两鬓斑白而十指漆黑。白与黑的对比，一是岁月，一是劳动。}
\notetext{4}{何所营：用来做什么？营，营生、维持。换钱不过为了一日温饱。}
\notetext{5}{愿天寒：希望天气更冷。身上穿得单薄，却盼天更冷——因为天冷炭才能卖好价钱。不是不怕冷，是更怕穷。}
\notetext{6}{辗冰辙：车轮碾过结了冰的车辙。辗（niǎn），碾过；辙（zhé），车轮印。}
\notetext{7}{黄衣使者白衫儿：穿黄衣的宦官和穿白衫的随从。黄衣，宦官服饰。}
\notetext{8}{口称敕：嘴里说着皇帝的命令。敕（chì），皇帝的诏令。}
\notetext{9}{回车叱牛牵向北：调转车头，呵斥着牛往北（皇宫方向）走。叱（chì），大声呵斥。}
\notetext{10}{宫使驱将惜不得：宫使把牛车赶走了，舍不得却无可奈何。"惜不得"三个字写尽了底层人的无力。}
\notetext{11}{半匹红纱一丈绫：半匹红纱和一丈绫罗。纱和绫都是丝织品，但价值远不及千余斤炭。}
\notetext{12}{充炭直：充当炭的价钱。直，同"值"。绑在牛头上充当炭钱——像施舍，实为掠夺。}
\end{annotations}

\begin{appreciation}
这首诗有一个副标题："苦宫市也"。宫市是唐代的一种制度：皇宫需要物资时，派宦官到民间采购，但往往只付很少的钱，甚至不付钱，形同掠夺。白居易写这首诗，是为了揭露宫市之害。

卖炭翁，伐薪烧炭南山中——开篇直入人物和场景。一个卖炭的老人，在终南山中砍柴烧炭。满面尘灰烟火色，两鬓苍苍十指黑——满脸是灰尘和烟火熏染的颜色，两鬓斑白，十个手指漆黑。这两句是肖像描写，也是一生劳作的缩影。苍苍与黑形成对比，白发是岁月的痕迹，黑指是劳动的印记。

卖炭得钱何所营？身上衣裳口中食——卖炭换来的钱用来做什么？不过是买身上的衣服、嘴里的食物。这是一个设问，答案简单得让人心酸。他的全部追求，不过是温饱而已。

可怜身上衣正单，心忧炭贱愿天寒——这是全诗最沉痛的一句。老人身上穿着单薄的衣服，却希望天气更冷一些，因为天冷了，炭才能卖个好价钱。冷与暖的矛盾，需求与现实的错位，把底层人的困境写得淋漓尽致。他不是不怕冷，是更怕穷。

夜来城外一尺雪，晓驾炭车辗冰辙——昨晚城外下了一尺厚的雪，天一亮就驾着炭车碾过结冰的车辙。辗（niǎn）是车轮滚过的意思。牛困人饥日已高，市南门外泥中歇——牛累了，人饿了，太阳已经很高了，才在长安城南门外的泥地里歇息。这一路的艰辛，只为了把炭卖到城里。

翩翩两骑来是谁？黄衣使者白衫儿——两个骑马的人轻快地来了，一个是穿黄衣的宦官，一个是穿白衣的随从。手把文书口称敕，回车叱牛牵向北——手里拿着公文，嘴里说着皇帝的命令，调转车头，呵斥着牛，把炭车拉向皇宫方向。敕（chì）是皇帝的诏令。叱（chì）是大声呵斥。两个同音字，一个是权力的象征，一个是权力的行使。

一车炭，千余斤，宫使驱将惜不得——一车炭，有一千多斤，宫里的使者赶走了牛车，老人舍不得却无可奈何。半匹红纱一丈绫，系向牛头充炭直——半匹红纱、一丈绫罗，绑在牛头上充当炭的价钱。绫（líng）是一种丝织品。千余斤炭，值多少钱？半匹红纱一丈绫，值多少钱？白居易没有说，但读者心里清楚。

这首诗的力量在于白描。白居易没有发表议论，只是把事情原原本本地记下来，愤怒和同情都在事实里。卖炭翁的悲苦，不在呼天抢地，而在那个"惜不得"——舍不得，但不得不舍。这种无力感，比任何控诉都更让人心碎。
\end{appreciation}
